%% One-sorted algebraic theories and the equivalence with equational algebras.

\subsection{Lawvere theories and equational algebras}
\label{sec:p10-lawvere-theories}

Throughout the Lawvere-theory development, signatures, equation sets,
algebra carriers, and generator sets belong to the fixed universe
used by \(\mathsf{Set}\) in Section~\ref{sec:p10-categorical-toolkit}.
The same constructions may be relativized to any larger universe
satisfying the size hypotheses of
Convention~\ref{def:p10-categorical-size}, when such a universe is fixed.

A finitary algebraic presentation can be recorded by its operations
in every finite list of variables and by the ways those operations
compose under substitution.  A Lawvere theory packages these data
as a category.  The number of variables is part of an arrow's type.

Lawvere's 1963 thesis uses the coproduct convention: the structural
functor goes from a category of finite sets to a theory category,
an \(n\)-ary operation goes from the generator to its \(n\)-fold
coproduct, and a model is a product-preserving functor from the
opposite theory to \(\mathsf{Set}\).
This is the presentation in Chapter~II, Section~1, pp.~61--63,
and Chapter~III, Section~1, pp.~74--75, of the authorized reprint
\cite{Lawvere2004FunctorialReprint}.
Passing to the opposite theory gives the product convention used
in Lawvere's separate 1968 article, Section~1, pp.~109--110 of that
same reprint, and in the Adamek--Rosicky--Vitale treatment below.
The two arrow conventions are related by this explicit reversal.

\begin{hbdefinition}{def:p10-lawvere-theory}{A one-sorted Lawvere theory}
Let \(\mathbb F\) be the category of finite ordinals
\(n=\{0,\ldots,n-1\}\) and all functions between them.
The empty ordinal is \(0\).
A one-sorted Lawvere theory consists of a small category \(T\)
whose objects are these natural numbers, and an identity-on-objects
functor
\[
                         j:\mathbb F^{\mathrm{op}}\longrightarrow T
\]
preserving finite products.
For the selector \(s_i:1\to n\), \(s_i(0)=i\), put
\(p_i^n=j(s_i^{\mathrm{op}}):n\to1\).
These projections exhibit \(n\) as the product of \(n\) copies
of the distinguished object \(1\).
In particular \(p_0^1=1_1\).
The object \(0\) is terminal; the distinguished generator \(1\)
is not by its name the terminal object.

For each family \(f_i:k\to1\), write
\(\langle f_i\rangle_{i<n}:k\to n\) for its unique product tuple.
For a function \(u:m\to n\) in \(\mathbb F\),
\[
                  j(u^{\mathrm{op}})
                     =\langle p_{u(i)}^n\rangle_{i<m}:n\to m.
\]
Thus structural maps select, repeat, or discard variables.
The functor \(j\) is not required to be faithful.
An \(n\)-ary operation is an arrow \(n\to1\); an arrow \(n\to m\)
is an \(m\)-tuple of \(n\)-ary operations.
A nullary operation is an arrow \(0\to1\).
\index{Lawvere theory}\index{algebraic theory!one-sorted}
\index{operation!as an arrow}\index{nullary operation}
\end{hbdefinition}

This is Definition~11.3 and Remarks~11.2--11.6, pp.~97--99,
of the author draft dated 1 February 2010
\cite{AdamekRosickyVitale2010Draft}.  The draft is a separately
identified version of the later published book.

\begin{hbdefinition}{def:p10-lawvere-presentation}{Equational presentations with specified finite contexts}
Fix a small finitary signature \(\Sigma\), allowing operation symbols
of arity zero.  Put \(X_n=(x_0,\ldots,x_{n-1})\), and let
\(T_\Sigma(n)=T_\Sigma(X_n)\) be its set of terms.
Here the term set is permitted to be empty: if \(n=0\) and there
are no nullary symbols, \(T_\Sigma(0)=\varnothing\).
This extends the carrier convention for the present construction;
Definition~\ref{def:p10-term-algebra} states the earlier nonempty
convention.

An equational presentation is a set \(E\) of triples
\((k,s,t)\), where \(s,t\in T_\Sigma(k)\).
The displayed context \(X_k\) records the variables over which
the identity is universally asserted.
A \(\Sigma\)-algebra with arbitrary carrier is a set \(A\) with
operations \(\sigma^A:A^k\to A\) of the declared arities.
The power \(A^k\) consists of functions from \(k\) to \(A\),
also written as \(k\)-tuples, and \(A^0\) is always a singleton.
The algebra satisfies \((k,s,t)\) when
\[
                     s^A(a)=t^A(a)\qquad\text{for every }a\in A^k.
\]
Evaluation uses the recursion of Definition~\ref{def:p10-term-algebra}.
An empty carrier admits all positive-arity operations uniquely,
but no nullary operation.
Write \(\operatorname{Alg}_{\mathrm{all}}(\Sigma,E)\) for these
models and their operation-preserving functions, and
\(\operatorname{Alg}_{\mathrm{ne}}(\Sigma,E)\) for the full
subcategory on nonempty carriers.

For every \(n\geq0\), let \(\equiv_E^n\) be the least family of
equivalence relations on \(T_\Sigma(n)\) that is compatible with
all operations, contains \(s\equiv_E^k t\) for \((k,s,t)\in E\),
and satisfies
\[
 s\equiv_E^k t
 \ \Longrightarrow\
 s[r_0/x_0,\ldots,r_{k-1}/x_{k-1}]
       \equiv_E^n
 t[r_0/x_0,\ldots,r_{k-1}/x_{k-1}]
       \quad(r_i\in T_\Sigma(n)).
\]
All substitutions are simultaneous.
The least family is the intersection of all families with these
properties; the family of full relations supplies one candidate.
It records the contextual form of equational deduction, including
instances of axioms written in other finite contexts.
\index{equational presentation}\index{algebra!empty carrier}
\index{equational consequence!finite contexts}
\end{hbdefinition}

\begin{hbtheorem}{thm:p10-lawvere-term-theory}{proved=annex:p10.lawvere.1}{The Lawvere theory of an equational presentation}
Define
\[
 T_{\Sigma,E}(n,m)=
             \bigl(T_\Sigma(n)/{\equiv_E^n}\bigr)^m.
\]
The identity of \(n\) is \(([x_i])_{i<n}\), and composition
substitutes the first arrow's tuple into each component of the
second arrow.
These data form a Lawvere theory.
Its structural functor sends \(u:m\to n\) to
\(([x_{u(i)}])_{i<m}:n\to m\).
Its terminal object is \(0\), and its binary product of \(n,m\)
is \(n+m\), with the two variable blocks as projections.
These assertions include empty source or target contexts.

For an algebra \(A\models E\), the formula
\[
 \begin{aligned}
 M_A(n)&=A^n,\\
 M_A(([t_i])_{i<m})(a)&=(t_i^A(a))_{i<m}
 \end{aligned}
\]
defines a finite-product-preserving functor
\(M_A:T_{\Sigma,E}\to\mathsf{Set}\).
An algebra homomorphism \(h:A\to B\) induces a natural
transformation \(M_h:M_A\Rightarrow M_B\) with components
\((M_h)_n=h^n\).
\end{hbtheorem}

The term theory and its evaluation are Notation~13.2 and
Lemma~13.3, p.~122, in
\cite{AdamekRosickyVitale2010Draft}; equations give the quotient
theory in Definition~13.10 and Theorem~13.11, p.~125.
The explicit context-category construction and its product
property are also given by
\cite[Section~4.2, pp.~30--35, especially
Proposition~4.2.1]{Pitts1995Draft}.

\begin{hbdefinition}{def:p10-lawvere-model}{Models and their underlying sets}
For a Lawvere theory \(T\), let
\[
                \operatorname{Mod}(T)
                    =\operatorname{FP}(T,\mathsf{Set})
\]
be the category of finite-product-preserving functors, with all
natural transformations as arrows.
The underlying-set functor is evaluation at the distinguished
generator:
\[
                 U(M)=M(1),\qquad U(\alpha)=\alpha_1.
\]
For every model the canonical comparison
\[
 c_n^M:M(n)\longrightarrow M(1)^n,\qquad
 c_n^M(z)=(M(p_i^n)(z))_{i<n},
\]
is a bijection.  For \(n=0\), this means that \(M(0)\) is
a singleton.  For \(n=1\), the comparison sends \(a\) to the
singleton tuple \((a)\).
Products are preserved through these comparisons; no strict
identity \(M(n)=M(1)^n\) is part of the definition.
\index{functorial semantics}\index{model!of an algebraic theory}
\index{underlying-set functor}
\end{hbdefinition}

\begin{hbtheorem}{thm:p10-lawvere-models}{proved=annex:p10.lawvere.2}{Equational algebras and functorial models}
Term evaluation gives an equivalence
\[
 \operatorname{Alg}_{\mathrm{all}}(\Sigma,E)
           \simeq\operatorname{Mod}(T_{\Sigma,E}).
\]
For the reverse construction, put \(A_M=M(1)\).
For a symbol \(\sigma\) of arity \(k\), let
\(\bar\sigma:k\to1\) be the one-component term arrow
\(([\sigma(x_0,\ldots,x_{k-1})])\), and define
\[
                     \sigma^{A_M}
                         =M(\bar\sigma)\circ(c_k^M)^{-1}.
\]
Then \(A_M\models E\), and a natural transformation
\(\alpha:M\Rightarrow N\) gives the algebra homomorphism
\(\alpha_1:A_M\to A_N\).
Every component of \(\alpha\) is determined by \(\alpha_1\):
\[
                     c_n^N\circ\alpha_n
                        =(\alpha_1)^n\circ c_n^M.
\]
The equivalence has explicit natural comparison isomorphisms.
On an algebra \(A\), the comparison with \(A_{M_A}=A^1\)
is \(a\mapsto(a)\); on a model \(M\), it is the family \(c_n^M\).

The equivalence restricts to
\[
 \operatorname{Alg}_{\mathrm{ne}}(\Sigma,E)
 \simeq
 \{M\in\operatorname{Mod}(T_{\Sigma,E}):M(1)\ne\varnothing\},
\]
where the right side is a full subcategory.
It therefore recovers the manuscript's earlier nonempty-algebra
convention without excluding empty set-valued models from the
unrestricted theorem.
\end{hbtheorem}

This functorial form, including all homomorphisms and natural
transformations, is
\cite[Theorem~13.11, p.~125; Proposition~11.8, p.~99]
{AdamekRosickyVitale2010Draft}
and \cite[Corollary~4.2.6(ii), p.~35]{Pitts1995Draft}.
Lawvere gives the original transformation-to-homomorphism
argument in Chapter~III, Section~1, p.~74 of the 1963 thesis
\cite{Lawvere2004FunctorialReprint}.

\begin{hbexample}{def:p10-lawvere-structural-collapse}{A structural functor can identify projections}
Take the empty signature and the equation
\((2,x_0,x_1)\).
The corresponding algebras have at most one element when empty
carriers are allowed.
In their term theory the two different selectors \(1\to2\) in
\(\mathbb F\) become the same projection \(2\to1\), because
\([x_0]=[x_1]\).
Hence the structural functor need not be faithful.
The object \(0\) still has no arrow to \(1\), since there are
no terms in the empty context.  This will distinguish its empty
initial model from the nonempty singleton models.
\end{hbexample}
